\documentclass[a4paper,11pt]{article}
\usepackage[T2A]{fontenc}
\usepackage[utf8]{inputenc}
\usepackage[russian]{babel}
\usepackage{amsmath,amssymb,mathtools}
\renewcommand{\familydefault}{\sfdefault}
\usepackage{geometry}
\geometry{left=21mm,right=21mm,top=20mm,bottom=21mm}
\usepackage{microtype}
\usepackage{array,booktabs,tabularx}
\usepackage{enumitem}
\usepackage{xcolor}
\usepackage{hyperref}
\usepackage{tikz}
\usepackage{fancyhdr}
\usepackage{parskip}
\definecolor{accent}{HTML}{16626A}
\definecolor{deep}{HTML}{26363A}
\definecolor{muted}{HTML}{58686D}
\definecolor{error}{HTML}{B02B35}
\hypersetup{colorlinks=true,linkcolor=accent,urlcolor=accent}
\setlength{\parindent}{0pt}
\setlength{\parskip}{4pt}
\setlist[itemize]{leftmargin=*,topsep=3pt,itemsep=2pt}
\pagestyle{fancy}
\fancyhf{}
\fancyfoot[C]{\footnotesize\color{muted}Лёвин Артём Александрович эксклюзивно для Ксении Клыковой}
\fancyfoot[R]{\footnotesize\color{muted}\thepage}
\renewcommand{\headrulewidth}{0pt}
\renewcommand{\footrulewidth}{0pt}
\newcommand{\blocktitle}[1]{\par\noindent{\bfseries\color{deep}#1}\par\nobreak}
\newcommand{\mathanswer}[1]{\textcolor{accent}{\textbf{Итог:}} #1}
\begin{document}
\thispagestyle{empty}
\vspace*{36mm}
\begin{center}
{\fontsize{13}{18}\selectfont\color{muted}МАТЕМАТИКА}\par\vspace{19mm}
{\fontsize{31}{38}\selectfont\bfseries\color{deep}Ревью домашней работы}\par\vspace{9mm}
{\fontsize{17}{23}\selectfont Банковские вклады и сложные проценты}\par
\vspace{28mm}
{\fontsize{17}{23}\selectfont Ксения Клыкова}\par
\vspace{11mm}
{\fontsize{13}{20}\selectfont 29 октября 2026 г.}\par
\vfill
{\fontsize{12}{19}\selectfont Лёвин Артём Александрович}\par
{\fontsize{11}{17}\selectfont\color{muted}Эксклюзивно для Ксении Клыковой}\par\vspace{11mm}
\begin{tikzpicture}[x=1cm,y=1cm]
\draw[accent!55,line width=0.8pt] (-5,0) .. controls (-2.7,0.55) and (-1.7,-0.4) .. (0,0) .. controls (1.8,0.45) and (3.2,-0.5) .. (5,0);
\end{tikzpicture}
\end{center}
\newpage
\begin{center}
{\LARGE\bfseries\color{deep}Сводная таблица ошибок}
\end{center}
\vspace{8mm}
В таблицу включены задания, в которых требуется уточнить или исправить ход решения. Ссылки в первом столбце ведут к подробному разбору.

\vspace{5mm}
\renewcommand{\arraystretch}{1.5}
\small
\begin{tabularx}{\linewidth}{@{}>{\raggedright\arraybackslash}p{12mm}>{\raggedright\arraybackslash}p{25mm}>{\raggedright\arraybackslash}X>{\raggedright\arraybackslash}p{21mm}>{\raggedright\arraybackslash}p{21mm}@{}}
\toprule
№ & Тема & Суть ошибки & Ваш ответ & Верный ответ \\
\midrule
\hyperref[task:5]{5} & Снятие денег со вклада & В промежуточной записи знаки перед суммами снятия не согласуются с исходной формулой; окончательное число верно. & 121{,}55 тыс. руб. & 121{,}55 тыс. руб. \\
\addlinespace
\hyperref[task:6]{6} & Пополнение вклада & Верная модель роста, но арифметическая ошибка при подсчёте суммы слагаемых. & 125{,}82 тыс. руб. & 125{,}52 тыс. руб. \\
\bottomrule
\end{tabularx}
\normalsize
\vfill
{\small\color{muted}Все денежные суммы в разборах приводятся в тысячах рублей, если не указано иное.}
\newpage
\thispagestyle{empty}
\null\vfill
\begin{center}
{\fontsize{26}{33}\selectfont\bfseries\color{deep}Разбор ошибок}\par\vspace{6mm}
{\fontsize{13}{19}\selectfont\color{muted}Сложные проценты, пополнения и снятия}
\end{center}
\vfill\null
\newpage
\phantomsection\label{task:5}
{\LARGE\bfseries\color{deep}Задание №\,5}\par\vspace{3mm}
\small
\blocktitle{Текст задания}
Вклад 100 тыс. руб. размещён на три года под 10\% годовых. В начале второго и третьего годов снимают по 5 тыс. руб. Найдите конечную сумму в тысячах рублей.
\blocktitle{Ваше решение}
В Вашей работе верно выделены три года начисления процентов и оба снятия:
\[
B_3=100\cdot 1{,}1^3-5\bigl(1{,}1^2+1{,}1\bigr).
\]
В следующей строке перед числами \(6{,}05\) и \(5{,}5\) записаны знаки сложения, хотя исходная формула содержит вычитание. Окончательный результат в работе: \(121{,}55\) тыс. руб.
\blocktitle{Ваш ответ}
\(121{,}55\) тыс. руб. Это правильный числовой ответ.
\blocktitle{Где именно возникает ошибка}
Последний правильный шаг --- формула для суммы после двух снятий. Первый ошибочный шаг --- переход от знака минус перед вкладом \(5(1{,}1^2+1{,}1)\) к знакам плюс перед отдельными слагаемыми.
\blocktitle{Почему это неверно}
Если из числа вычитают сумму, нужно вычесть каждое её слагаемое:
\[
a-(b+c)=a-b-c.
\]
Снятие денег уменьшает итоговый вклад. Запись с плюсами дала бы \(133{,}1+6{,}05+5{,}5=144{,}65\) тыс. руб., что больше суммы без снятий. Такой результат противоречил бы смыслу задачи.
\blocktitle{Ключевая идея}
\textcolor{accent}{\textbf{Ключевая идея:}} каждое снятие сначала уменьшает вклад, затем оставшаяся сумма продолжает приносить проценты. Снятие в начале второго года влияет на два года, а снятие в начале третьего --- на один.
\blocktitle{Исправление от последнего правильного шага}
\[
\begin{aligned}
B_3&=100\cdot1{,}1^3-5\bigl(1{,}1^2+1{,}1\bigr)\\
&=133{,}1-6{,}05-5{,}5.
\end{aligned}
\]
\blocktitle{Полное правильное решение}
Первоначальные 100 тыс. руб. начисляют проценты трижды, поэтому к концу срока без снятий они превратились бы в
\[
100\cdot1{,}1^3=100\cdot1{,}331=133{,}1.
\]
Первые 5 тыс. руб. снимают в начале второго года. До конца срока остаются два начисления процентов, поэтому их влияние на итог: \(5\cdot1{,}1^2=6{,}05\).

Вторые 5 тыс. руб. снимают в начале третьего года. До конца срока остаётся одно начисление: \(5\cdot1{,}1=5{,}5\).

Вычитаем оба эффекта снятия:
\[
B_3=133{,}1-6{,}05-5{,}5=\boxed{121{,}55}.
\]
\blocktitle{Проверка}
По годам: после первого года \(110\), затем снимаем \(5\), получаем \(105\); после второго года \(105\cdot1{,}1=115{,}5\), затем снимаем \(5\), получаем \(110{,}5\); после третьего года \(110{,}5\cdot1{,}1=121{,}55\). Результаты совпадают.
\blocktitle{Итог}
\mathanswer{\(121{,}55\) тыс. руб. Число найдено верно; исправить требуется запись промежуточных знаков.}
\blocktitle{Задача на закрепление}
Вклад 150 тыс. руб. размещён на три года под 10\% годовых. В начале второго и третьего годов снимают по 6 тыс. руб. Сколько тысяч рублей останется на счёте к концу третьего года?
\newpage
\phantomsection\label{task:6}
{\LARGE\bfseries\color{deep}Задание №\,6}\par\vspace{3mm}
\small
\blocktitle{Текст задания}
Вклад 50 тыс. руб. размещён на четыре года под 20\% годовых. В начале второго, третьего и четвёртого годов добавляют по 5 тыс. руб. Найдите сумму к концу четвёртого года в тысячах рублей.
\blocktitle{Ваше решение}
Вы составили правильную формулу, учитывая, сколько времени каждое пополнение приносит проценты:
\[
B_4=50\cdot1{,}2^4+5\cdot1{,}2^3+5\cdot1{,}2^2+5\cdot1{,}2.
\]
\blocktitle{Ваш ответ}
\(125{,}82\) тыс. руб.
\blocktitle{Ошибка и точное место}
\textcolor{error}{\textbf{Ошибка:}} арифметическая ошибка при вычислении записанной суммы. Последний правильный шаг --- формула для \(B_4\); первый неверный --- получение числового значения \(125{,}82\) по этой формуле.
\blocktitle{Почему это неверно}
Каждая степень отвечает отдельному сроку начисления процентов. Их коэффициенты и показатели Вы выбрали верно. Ошибка возникает только при вычислении одного или нескольких слагаемых либо при их сложении. Чтобы её обнаружить, следует посчитать каждое слагаемое отдельно и только затем сложить результаты.
\blocktitle{Ключевая идея}
\textcolor{accent}{\textbf{Ключевая идея:}} первоначальный вклад растёт четыре года; пополнение в начале второго года --- три года; в начале третьего --- два года; в начале четвёртого --- один год.
\blocktitle{Исправление от последнего правильного шага}
\[
\begin{aligned}
50\cdot1{,}2^4 &=50\cdot2{,}0736=103{,}68,\\
5\cdot1{,}2^3 &=5\cdot1{,}728=8{,}64,\\
5\cdot1{,}2^2 &=5\cdot1{,}44=7{,}20,\\
5\cdot1{,}2 &=6{,}00.
\end{aligned}
\]
\blocktitle{Полное правильное решение}
Сумма к концу четвёртого года равна сумме будущих стоимостей первоначальных денег и каждого пополнения:
\[
\begin{aligned}
B_4&=50\cdot1{,}2^4+5\cdot1{,}2^3+5\cdot1{,}2^2+5\cdot1{,}2\\
&=103{,}68+8{,}64+7{,}20+6{,}00\\
&=\boxed{125{,}52}.
\end{aligned}
\]
\blocktitle{Проверка}
Посчитаем сумму по годам. После первого года: \(50\cdot1{,}2=60\). Пополнение: \(60+5=65\). После второго года: \(65\cdot1{,}2=78\). Пополнение: \(78+5=83\). После третьего года: \(83\cdot1{,}2=99{,}6\). Пополнение: \(99{,}6+5=104{,}6\). После четвёртого года: \(104{,}6\cdot1{,}2=125{,}52\).
\blocktitle{Итог}
\mathanswer{\(125{,}52\) тыс. руб., или \(125\,520\) руб.}
\blocktitle{Задача на закрепление}
Вклад 40 тыс. руб. размещён на четыре года под 15\% годовых. В начале второго, третьего и четвёртого годов добавляют по 4 тыс. руб. Найдите сумму к концу четвёртого года в тысячах рублей.
\end{document}
